\scanpage{282}
\mspage{12}

\uline{Loc} (5) \inserted{(Axiome de \uline{connexité}) ou de \uline{non-trivialité})} Soit $\unsure{f}\in W$, alors $\pi_0(f)$ est bijectif.

Cette \deleted{axiome} propriété de $W$ est d'une \unsure{utilité} \unsure{manifeste|constante} en pratique, quand on travaille avec $W$ au lieu de se borner à construire des $W$ et faire joujou avec les axiomes et leur stabilité par passage à des $\underline{\operatorname{Hom}}(I,M)$ (où ici $M=\mathrm{Cat}$). D'autre part, on peut aussi la regarder comme un axiome de \uline{non-trivialité} pour $\mathrm{Hot}_W$. Je rappelle à ce sujet : \deleted{\illegible{1}}

\uline{Proposition} 2.7. Soit $W$ un localiseur fondamental. Conditions équivalentes

a) $W$ satisfait l'axiome de connexité \uline{Loc} 5

b) Le foncteur $\pi_0:\mathrm{Cat}\unsure{{-}|\longrightarrow}\mathrm{Ens}$ se factorise par un foncteur

\diagram{\begin{tikzcd}[ampersand replacement=\&,column sep=large]
\mathrm{Cat} \arrow[r,dotted] \arrow[rr,bend right=24,dashed,"{\pi_0}"'] \& \mathrm{Hot}_W \arrow[r,"{\pi_0}"] \& \mathrm{Ens},
\end{tikzcd}}

% slanted
\margin{\uline{NB} Si $W$ est fort saturé (condition $\operatorname{Loc}(1\,\mathrm{ter})$), alors les conditions ci-contre (i.e. $\operatorname{Loc}(5)$) équivalent à ceci : \deleted{\unsure{b|a}}\unsure{d)} $W$ n'est égal ni à $\mathrm{Fl}(\mathrm{Cat})$ tout entier, ni à $\mathrm{Fl}(\mathrm{Cat}^{*})\cup\{\mathrm{id}_{\varnothing}\}$ (i.e. $\mathrm{Fl}(\mathrm{Cat})\setminus W$ n'est ni vide, ni égal à l'ensemble des flèches $\varnothing\longrightarrow X$ où $X\ne\varnothing$). Si $W$ satisfait à $\operatorname{Loc}(\unsure{7\,D_{W_{\omega}}|7\,D_{W\omega}|7\,D_{Ww}})$, alors $W$ \deleted{ne} satisfait $\operatorname{Loc}(5)$ sss $W\ne\mathrm{Fl}(\mathrm{Cat})$ (cf XII th. 2, p. 20)}

c) La catégorie $\mathrm{Hot}_W^{*}=\mathrm{Hot}_W\setminus\{\varnothing\}$ n'est équivalente à la catégorie ponctuelle

La catégorie $\mathrm{Hot}_W$ n'est pas équivalente \deleted{\unsure{à}} ni à $\unsure{\mathbb{\Delta}|\Delta}^{1}$.

Pour la démonstration et pour d'autres formes équivalentes de l'axiome de \unsure{connexité},

% archive sheet number: 281
